\documentclass[11pt,a4paper]{article}

% ---- packages ----
\usepackage[margin=1in]{geometry}
\usepackage{amsmath}
\usepackage{booktabs}
\usepackage{array}
\usepackage{enumitem}
\usepackage{xcolor}
\usepackage{listings}
\usepackage[colorlinks=true,linkcolor=blue,urlcolor=blue]{hyperref}

% ---- listings style ----
\definecolor{codebg}{rgb}{0.96,0.96,0.96}
\lstset{
  basicstyle=\ttfamily\small,
  backgroundcolor=\color{codebg},
  breaklines=true,
  frame=single,
  framesep=4pt,
  columns=fullflexible,
  keepspaces=true,
  showstringspaces=false,
}

\title{\textbf{Preserved-Identity}\\[4pt]
\large Rearranging an Image's Own Pixels to Take the Form of Another}
\author{Yash Yadav}
\date{June 2025}

\begin{document}
\maketitle

\begin{abstract}
\noindent
Given a \emph{source} image (a ``bag of pixels'') and a \emph{target} image (a form to
reproduce), this project produces an output that \textbf{looks like the target but is built
entirely from the source's own pixels}. The output is a strict \emph{permutation} of the
source --- every pixel used exactly once --- so it preserves the source's exact colour
histogram, contrast, and brightness while taking on the target's spatial structure. We
formulate this as an optimal-transport assignment problem, solve it at scale with a recursive
median-split approximation, and apply it in a source-anchored portrait pipeline that morphs a
photo's subject onto a target face while keeping the original image intact.
\end{abstract}

% =====================================================================
\section{Objective}
Given a \textbf{source} image (e.g.\ a photo of a person outdoors) and a \textbf{target}
image (e.g.\ a portrait), produce an output that resembles the target but is composed
\textbf{entirely from the source's own pixels}.

The defining constraint is that the output is a \textbf{permutation} of the source: every
source pixel is used exactly once, none invented and none dropped. Consequently the output
keeps the source's exact \textbf{colour histogram, contrast, and brightness} --- it is
literally the same pixels, relocated in space. Dark source pixels migrate to where the target
is dark (hair, shadows), bright pixels to highlights, and so on.

This is \textbf{not} a photomosaic (which tiles a picture out of many small \emph{separate}
images). It is a single-image \textbf{pixel-rearrangement / optimal-transport} problem.

% =====================================================================
\section{Problem Formulation}
Let the source contribute $N$ pixels and the target define $N$ positions. We seek a bijection
$\sigma$ (a one-to-one assignment) from source pixels to target positions that minimises the
total perceptual colour distance:
\begin{equation}
\min_{\sigma \in \mathcal{S}_N}
\sum_{p=1}^{N} \operatorname{dist}\!\big(\, c^{\text{src}}_{\sigma(p)} \,,\, c^{\text{tgt}}_{p} \,\big),
\end{equation}
where $\mathcal{S}_N$ is the set of all permutations of $N$ elements, $c^{\text{src}}$ and
$c^{\text{tgt}}$ are source and target colours, and $\operatorname{dist}(\cdot,\cdot)$ is
Euclidean distance in \textbf{CIELAB} colour space (perceptually uniform, closer to how the
eye judges colour than raw RGB). Solving this exactly is the \textbf{linear assignment
problem}.

% =====================================================================
\section{The Matching Engine (\texttt{src/pixel\_rearrange.py})}
Three interchangeable solvers trade quality for scale:

\begin{center}
\begin{tabular}{@{}llll@{}}
\toprule
\textbf{Method} & \textbf{Technique} & \textbf{Quality} & \textbf{Scale} \\
\midrule
\texttt{exact} & Hungarian algorithm & Optimal & $\le 64{\times}64$ \ ($O(N^3)$) \\
\texttt{recursive}\, \emph{(default)} & Median-split optimal transport & $\sim$few \% from optimal & Megapixels ($\sim$2\,s @ $1024^2$) \\
\texttt{sorted} & Lightness sort & Coarse & Instant \\
\bottomrule
\end{tabular}
\end{center}

\paragraph{Recursive median-split OT} is the workhorse that made the project scale. It
recursively partitions both pixel clouds along their highest-variance colour axis, splitting
each at the median so the halves keep equal counts, and solves the small blocks
($\le \texttt{leaf}$ pixels) exactly with the Hungarian algorithm. This preserves a true global
bijection while running in $\sim O(N \log N)$ plus cheap leaf solves --- turning an
intractable $O(N^3)$ assignment into one that handles full photographs.

\paragraph{Engine extras.}
\begin{itemize}[leftmargin=1.4em]
  \item \textbf{Recolour knob} (\texttt{-{}-recolor}): optionally nudge rearranged pixels
  toward the true target colours. \texttt{luma} mode adopts the target's lightness/contrast
  while keeping the source's chroma; \texttt{blend} is a plain RGB blend. $0$ = pure shuffle
  (histogram perfectly preserved).
  \item \textbf{Priority mode}: the most important target positions (face/edges) claim the
  globally best-fitting source pixels first, while remaining a strict permutation.
  \item \textbf{Aspect handling}: the output follows the target's native aspect ratio.
\end{itemize}

% =====================================================================
\section{Source-Anchored Portrait Pipeline (\texttt{src/portrait.py})}
The headline application keeps the source image intact (size, palette, contrast) and morphs
\textbf{only its subject} onto the target's form, on a clean background --- a
\textbf{budget-driven, gap-free} rearrangement:

\begin{enumerate}[leftmargin=1.6em]
  \item \textbf{Partition the source by count.} Segment into the dominant background (e.g.\
  sky) and everything else (the ``object pool''), giving $N_{\text{obj}}$ object pixels and
  $N_{\text{bg}}$ background pixels.
  \item \textbf{Size the subject to the budget.} Take the target's subject silhouette
  (background removed, auto-oriented upright) and \textbf{scale it so its area equals
  $N_{\text{obj}}$}, centred in the frame. The remaining border then equals $N_{\text{bg}}$
  exactly --- counts match on both sides.
  \item \textbf{Overshoot + erode.} The subject is scaled to slightly over-fill, then eroded
  to the exact budget. Growing a region (dilation) produces blobby halos, whereas trimming an
  oversized subject (erosion) stays clean.
  \item \textbf{Render the subject.} The $N_{\text{obj}}$ object pixels are rearranged onto the
  subject guide with \textbf{tonal cluster sub-matching} (dark $\rightarrow$ hair/shadow, mid
  $\rightarrow$ skin) for a face-like read.
  \item \textbf{Smooth background.} The $N_{\text{bg}}$ background pixels fill every border
  position, each taking its \textbf{nearest original background colour} --- reproducing the
  sky's gradient and blending the vacated centre seamlessly.
\end{enumerate}

\noindent
\textbf{Result:} every source pixel used once, every output position filled --- no
empty/jagged gaps, no halos, identical histogram/contrast/brightness to the source, subject
cleanly centred.

% =====================================================================
\section{Segmentation (\texttt{src/segment.py})}
Semantic, \textbf{not} colour-based --- which lets a blue book held against a blue wall stay
foreground while the wall becomes background:

\begin{itemize}[leftmargin=1.4em]
  \item \textbf{Mask R-CNN} (COCO instances) --- unions every confidently detected object. A
  \texttt{person\_only} flag restricts to the \texttt{person} class for the \emph{target}
  subject guide, so nearby furniture (a chair behind the head) is never merged into the
  silhouette. The source still unions all objects to protect held items.
  \item \textbf{DeepLabV3} (person class) --- fallback when Mask R-CNN is unavailable.
  \item \textbf{GrabCut} (OpenCV) --- final offline fallback.
  \item \textbf{Face-aware auto-orientation} (MTCNN) --- detects the most confident face across
  the four right-angle rotations and rotates the target upright before fitting.
  \item \textbf{Largest-connected-component} filtering removes stray detached mask blobs.
\end{itemize}

% =====================================================================
\section{Module Map}
\begin{center}
\begin{tabular}{@{}>{\ttfamily}l p{0.62\textwidth}@{}}
\toprule
\textnormal{\textbf{File}} & \textbf{Role} \\
\midrule
src/pixel\_rearrange.py & Optimal-transport assignment engine (exact / recursive / sorted), recolour, priority. \\
src/portrait.py & Source-anchored portrait pipeline (the main product). \\
src/segment.py & Mask R-CNN / DeepLabV3 / GrabCut masks, face detection \& auto-orient. \\
src/saliency.py & Target importance map (centre bias + edge energy + face box) for \texttt{-{}-priority}. \\
src/cli.py & CLI for plain whole-image rearrangement. \\
scripts/demo\_rearrange.py & Synthetic tree $\rightarrow$ face proof-of-concept. \\
\bottomrule
\end{tabular}
\end{center}

% =====================================================================
\section{How to Run}
\textbf{Source-anchored portrait (main):}
\begin{lstlisting}[language=bash]
python -m src.portrait \
    --source data/source.jpeg \
    --target data/target.jpg \
    --out results/portrait.png \
    --max-dim 1600 --seg maskrcnn
\end{lstlisting}

\noindent\textbf{Plain whole-image rearrangement:}
\begin{lstlisting}[language=bash]
python -m src.cli \
    --source data/source.jpg --target data/target.jpg \
    --out results/out.png --size 1024 --space lab --recolor 0.3
\end{lstlisting}

\noindent
Install: \texttt{pip install -r requirements.txt} (numpy, scipy, Pillow, scikit-image, plus
opencv-python / facenet-pytorch / torch / torchvision for the portrait pipeline).

% =====================================================================
\section{Choosing a Good Source (Matters Most)}
Because the output is a \emph{pure rearrangement}, its colours are \textbf{locked to the
source's palette} --- no algorithm can create colours the source lacks. A mostly-white source
can never produce dark hair; a dark source can never make a bright sky. For best results pick a
source whose brightness/colour distribution \textbf{overlaps the target's} (mid-tones for skin,
some dark pixels for hair/shadows, etc.).

% =====================================================================
\section{Development Journey}
\begin{enumerate}[leftmargin=1.6em]
  \item \textbf{Course correction.} The repository began as an identity-aware
  \emph{photomosaic} (face tiles + embeddings + LPIPS). That solved the wrong problem; it was
  scrapped in favour of single-image pixel rearrangement.
  \item \textbf{Proof of concept.} Exact Hungarian assignment in RGB on small synthetic images
  (tree $\rightarrow$ face) confirmed the histogram-preserving rearrangement worked.
  \item \textbf{Quality.} Switched matching to perceptual CIELAB --- visibly smoother, more
  faithful.
  \item \textbf{Scale.} Added the recursive median-split OT solver to reach megapixel images.
  \item \textbf{Portrait pipeline.} Built segmentation plus the source-anchored portrait so the
  subject takes the target's form while the source is preserved.
  \item \textbf{Final polish.} Replaced the frozen-background approach with the budget-driven
  partition (gap-free, smooth background, subject centred); fixed a blobby halo via
  overshoot-then-erode and largest-component filtering; and fixed background leakage into the
  subject by restricting the target mask to \texttt{person\_only}.
\end{enumerate}

% =====================================================================
\section{Limitations \& Future Work}
\begin{itemize}[leftmargin=1.4em]
  \item \textbf{Palette lock} is fundamental --- resemblance depends on source/target palette
  overlap.
  \item The recursive solver is an \emph{approximation} of optimal transport (a few \% from
  optimal); Sinkhorn could close the gap if needed.
  \item Segmentation quality bounds portrait quality; unusual poses or occlusions can still
  mis-segment. Soft-alpha matting would refine subject boundaries.
  \item Background smoothing assumes a single dominant background region (works well for
  sky/wall; busy backgrounds are less clean).
\end{itemize}

% =====================================================================
\section{Status}
\textbf{Complete.} The pipeline produces clean, gap-free source-anchored portraits: the
target's form rendered entirely from the source's own pixels, subject centred on a smooth
background, with the source's exact histogram, contrast, and brightness preserved.

\end{document}
